% chapters/09/praktikum.tex
\begin{praktikumbox}
Praktikum ini mengikuti Tugas Praktikum 07 pada RPS: NPC berpatroli di antara
titik-titik tertentu, beralih mengejar ketika pemain masuk jari-jari persepsi,
dan kembali berdiam ketika pemain keluar dari jangkauan.

\begin{langkah}
  \item \textbf{Jalankan programnya.} Ketik ulang \code{ch09_fsm.py} dari
        Daftar~\ref{lst:ch09-fsm} ke dalam folder \code{code}. Jalankan, lalu
        amati warna NPC yang berubah menurut keadaannya dan jalur patroli yang
        dilaluinya.

  \item \textbf{Amati setiap keadaan secara terpisah.} Berdiamlah di jauh
        sampai NPC menyelesaikan satu putaran patroli, lalu dekati sampai ia
        mengejar, lalu berdirilah di dekatnya sampai ia menyerang. Catat
        keadaan mana yang paling lama berlangsung.

  \item \textbf{Jalankan skenario pengujian.} Jalankan program dengan pilihan
        \code{--scenario}, lalu catat seluruh catatan perpindahan keadaan
        beserta nomor bingkainya.
\begin{kodekeluaran}
python ch09_fsm.py --frames 900 --scenario
\end{kodekeluaran}

  \item \textbf{Gambarlah diagram keadaan Anda sendiri.} Gambarlah diagram
        keadaan--transisi untuk NPC Anda, dengan empat keadaan dan syarat pada
        setiap perpindahannya. Bandingkan dengan
        Gambar~\ref{fig:diagram-keadaan}, lalu periksa apakah setiap syarat
        pada diagram Anda benar-benar ada di dalam kode.

  \item \textbf{Tambahkan satu keadaan baru.} Tambahkan keadaan menghindar,
        yang dijalani NPC ketika nyawanya tinggal sedikit. Tentukan syarat
        masuk dan syarat keluarnya, gambar pada diagram, lalu tambahkan pada
        kode beserta catatan perpindahannya.

  \item \textbf{Uji perpindahan yang tidak sah.} Tambahkan satu perpindahan
        yang tidak ada pada diagram Anda, misalnya dari menyerang langsung ke
        berpatroli, lalu jalankan skenario pengujian. Perhatikan laporan
        pemeriksaan keabsahan pada terminal, lalu kembalikan program.

  \item \textbf{Tulis catatan praktikum} yang memuat diagram keadaan Anda,
        catatan perpindahan beserta nomor bingkainya, dan hasil pemeriksaan
        keabsahan.

  \item \textbf{Simpan pekerjaan} ke repositori Git dengan pesan penyimpanan
        yang menerangkan isinya.
\end{langkah}
\end{praktikumbox}

\begin{table}[htbp]
\centering
\caption{Lembar kerja Praktikum 07}
\label{tab:lembar-praktikum-07}
{\setlength{\tabcolsep}{4pt}\footnotesize%
\begin{tabularx}{\textwidth}{@{}C{3.4cm} C{2.6cm} C{2.6cm} Y@{}}
\toprule
\headrow \textbf{Keadaan} & \textbf{Lama berlangsung} & \textbf{Jumlah perpindahan} & \textbf{Catatan} \\
\midrule
berdiam & & & \\[1.2em]
\addlinespace
berpatroli & & & \\[1.2em]
\addlinespace
mengejar & & & \\[1.2em]
\addlinespace
menyerang & & & \\[1.2em]
\addlinespace
menghindar (keadaan tambahan) & & & \\[1.2em]
\addlinespace
\multicolumn{4}{@{}p{\textwidth}@{}}{\footnotesize\itshape
Hasil pemeriksaan keabsahan perpindahan: \dotfill} \\[1.4em]
\addlinespace
\multicolumn{4}{@{}p{\textwidth}@{}}{\footnotesize\itshape
Perpindahan yang dilaporkan tidak sah pada langkah keenam: \dotfill} \\[1.4em]
\bottomrule
\end{tabularx}}
\end{table}

\begin{catatabox}
Langkah kelima dan keenam saling melengkapi. Pada langkah kelima Anda
menambahkan keadaan beserta perpindahannya, sedangkan pada langkah keenam Anda
menambahkan perpindahan yang tidak ada pada diagram. Pemeriksaan keabsahan akan
menangkap yang kedua dan tidak mengganggu yang pertama, asalkan diagramnya
diperbarui lebih dahulu.
\end{catatabox}
